@@M0@@ and the degree filtration on the second member follows readily from (V 6.6.4) and from the fact that the @@M1@@ define a \emph{Chern theory} of @@M2@@ in the sense of (V 6.13), when @@M3@@ is regarded as a @@M4@@-algebra augmented over the binomial ring @@M5@@). One thereby obtains maps
\sgaStableTarget{sga6.e10.eq.7-13-7}@@M6@@
by the formula
\sgaStableTarget{sga6.e10.eq.7-13-8}@@M7@@
compare (XIV 5.1 and 5.2). With this understood, the commutativity of diagrams (7.13.2) and (7.13.4) trivially implies that of the analogous diagrams
\sgaStableTarget{sga6.e10.eq.7-13-9}@@M8@@
\sgaStableTarget{sga6.e10.eq.7-13-10}@@M9@@
\sgaStableTarget{sga6.e10.rem.7-13-11}\textbf{Remark 7.13.11.} The commutativity of the last two diagrams, which can be expressed by saying that specialization of algebraic cycles is compatible with formation of the associated @@M10@@-adic cohomology classes, is not very satisfactory in the form stated, since we have had to replace the coefficients @@M11@@ by @@M12@@, that is, to neglect torsion. To obtain a statement that does not neglect torsion, it seems that one
